% !TEX TS-program = lualatex+se
% !TEX encoding = UTF-8
\documentclass[12pt, a4paper]{book}

% --- USTAWIENIA STRONY ---
\usepackage{geometry}
\geometry{
%	showframe=true,
%	showcrop=true,
	left=15mm,
	right=15mm,
	top=15mm,
	bottom=15mm,
	marginparwidth=0mm,
	marginparsep=0mm
	}
%\usepackage[letterpaper,showframe]{geometry}	
% usuwanie nagłówków/stopek z\,dokumentu
\usepackage[indent=0pt]{parskip} % ustawianie paragrafów
%\setstretch{1.5} % półtorej linii odległości między wierszami
\usepackage{titlesec} % ustawianie nagłówków
%\usepackage{fullpage} % redukcja marginesów
\usepackage{setspace} % odstępy między liniami
\usepackage{enumitem}% listy numerowane
% -----------------------------

%----------- style nagłówków ----------------
\newcommand{\headingone}{\scshape\Large}
\newcommand{\headingtwo}{\large}
\newcommand{\margintext}[1]{\textit{\small #1}}
% przesunięcia poziome
\newcommand{\smallwidth}{\hskip 1em}
\newcommand{\medwidth}{\hskip 3pt}
\newcommand{\largewidth}{\hskip 12pt}
% przesunięcia pionowe
\newcommand{\linespace}{\vskip 6pt}
\newcommand{\paragraphspace}{\vskip 10pt}
\newcommand{\segmentspace}{\vskip 16pt}
\newcommand{\sectionspace}{\vskip 28pt}
%------------------------------------------

% --- GREGORIO ---	
\usepackage{gregoriotex}
\gresetlinecolor{gregoriocolor}
\gresetheadercapture{annotation}{greannotation}{string}% odczytuje adnotacje nad inicjałem jako ciąg, co pozwala na jego renderowanie
\gresetheadercapture{commentary}{grecommentary}{string}% odczytuje komentarz jako ciąg, co pozwala na jego renderowanie
\grechangedim{commentaryraise}{0.3cm}{scalable} % podnosi komentarz wystarczająco ponad górną linię
\grechangestyle{annotation}{\color{black}\small}% format adnotacji
\grechangestyle{commentary}{\color{gregoriocolor}\small\itshape}% format komentarza
\grechangestyle{abovelinestext}{\color{gregoriocolor}\small\itshape}% format napisu nad liniami
\grechangestyle{initial}{\fontsize{36}{36}\selectfont} % zmniejsza wielkość inicjału do bardziej proporcjonalnego rozmiaru (=36pt)
\grechangedim{intersyllablespacenotes}{0.26cm}{scalable}% Minimalna odległość między nutami różnych sylab.

% umieszczenie zawartości składu ich unikalnego argumentu po lewej i prawej stronie każdej linii
\let\grelocalleftbox\localleftbox
\let\grelocalrightbox\localrightbox

%\newcommand{\separator}{\greseparator{3}{40}}
%\newfontface\vectis{VectisMonumentalW01-Regular}
%\def\greinitialformat#1{\fontsize{20}{20}\vectis #1}
% ------------------

%\usepackage[hyperindex=true]{hyperref}% przypisy dolne

%---- usun konflikt między paracol i luacolor ----
%-------- kod poprawiający konflikt między paracol i luacolor ----
\ifluatex
\usepackage{luacolor}
\directlua{
 resetboxcolor = function (head,color)
  for n in node.traverse(head) do
    if 
      node.has_attribute(n,oberdiek.luacolor.getattribute(),
        oberdiek.luacolor.getvalue(token.get_macro("default@color")))
    then 
      node.set_attribute(n,oberdiek.luacolor.getattribute(),color)      
    end
    if n.id == node.id("hlist") or n.id == node.id("vlist") then
     resetboxcolor (n.list,color)
    end
 end
 return head
 end 
}

\protected\def\luacolorresetboxcolor#1{%
  \directlua {local color= tex.getattribute(oberdiek.luacolor.getattribute())
              resetboxcolor (tex.getbox(\the#1),color)}}

\else
\protected\def\luacolorresetboxcolor#1{}
\fi
%------ koniec kodu --------------------------------------
%-----------------------------------------------

% --- SPIS TREŚCI -------
%\usepackage{imakeidx}
%\makeindex[title=Alfabetyczny spis program=texindy,options=-M styl.xdy]
%\usepackage{filecontents}
%\begin{filecontents*}{styl.xdy}
%;;; xindy style file
%(markup-locclass-list :open "\dotfill" :sep "")
%(define-attributes (( "textbf" "default" )) )
%(markup-locref   :attr  "textbf"     :open "\textbf{\hyperpage{" :close "}}")
%(markup-locref   :attr  "textit"     :open "\textit{\hyperpage{" :close "}}")
%(markup-locref   :attr  "textttt"     :open "\textttt{\hyperpage{" :close "}}")
%(markup-locref   :attr  "texttsc"     :open "\texttsc{\hyperpage{" :close "}}")
%(markup-locref   :attr  "default"     :open "\hyperpage{" :close "}")
%\end{filecontents*}
% -------------------------

% ---  Konfiguracja kolumn ---
\usepackage{paracol}
\setcolumnwidth{.49\textwidth, .49\textwidth}
\setlength{\columnseprule}{0.4pt}
\setlength{\columnsep}{.02\textwidth}
\colseprulecolor{red}% wymagany pakiet xcolor
\footnotelayout{m}
% ----------------------------

\makeatletter 
\patchcmd\pcol@hfil
 {\copy\pcol@tempboxa}
 {\luacolorresetboxcolor\pcol@tempboxa\copy\pcol@tempboxa}{}{\fail}
\makeatother
% ----------------------------

\footnotelayout{m}

%% --- UKŁAD DWUKOLUMNOWY Z TŁUMACZENIEM ---
%\usepackage{tcolorbox}
%\newtcolorbox{nutybox}{size=minimal, colback=white, colframe=red, arc=0mm, width=\linewidth, halign=left, valign=center}
%\newtcolorbox{tlumbox}{size=minimal, colback=white, colframe=red, arc=0mm, width=\linewidth, halign=left, valign=top, fontupper=\small\itshape} 
%\newcommand{\nuty}[1]{\begin{nutybox}#1 \end{nutybox}}
%\newcommand{\tlum}[1]{\begin{tlumbox}\tlumacz{#1} \end{tlumbox}}
%% ------------------------------------------------------

\usepackage{Tabbing}
\usepackage{longtable}
\usepackage[savepos]{zref}

% --- CZCIONKI DOKUMENTU ---
\usepackage[latin,polish]{babel}
\usepackage{fontspec}
%\usepackage{times}
%\usepackage[math]{anttor}
%\usepackage{tgtermes}
%\usepackage{mathptmx}
\usepackage{ebgaramond}
\usepackage{libertine}
\defaultfontfeatures{Ligatures=TeX}

%\setmainfont{TeX Gyre Termes}
%\setmainfont{TeX Gyre Heros}
%\setmainfont{Calluna Regular}
\setmainfont{EB Garamond}
%\setmainfont{Bodoni Moda}
%\newfontfamily{\sekfont}{Times New Roman}
\newfontfamily{\sekfont}{Libertinus Serif}
\newcommand\sekcjafont{\fontsize{16}{16}\selectfont \sekfont}

%----- Strona tytułowa -------------
\newfontfamily{\tytfont}{JSL Ancient}
\newcommand\tytulfont[1]{\fontsize{36}{30}\selectfont \tytfont#1\par}
\newcommand\podtytulfont[1]{\fontsize{18}{10}\selectfont \tytfont#1\par}

%\setmainfont[<font features>]{<name of font>} % sets the roman font
%\setsansfont[<font features>]{<name of font>} % sets the sans font
%\setmonofont[<font features>]{<name of font>} % sets the monospace font

% --- KONIEC CZCIONKI DOKUMENTU ---

% --- GRAFIKA I KOLORY ---
\usepackage{xcolor} % kolory
\usepackage{graphicx} % umożliwia import obrazów
\graphicspath{{img/}} % katalog obrazów
\usepackage{pgfornament}% ornamenty graficzne
% --- KONIEC GRAFIKA I KOLORY -----------

% --- KOMENDY WŁASNE ---
\newcommand{\red}[1]{\textcolor{red}{#1}}
\newcommand{\black}[1]{\textcolor{black}{#1}}
% --------------------------------
\newcommand{\tlumacz}[1]{\textit{#1}}
\newcommand{\inicjal}[2]{\textcolor{red}{\LARGE\textbf{#1}}{#2}}
%--------------------------------
\titlespacing*{\chapter}{0pt}{0pt}{5pt}	
\titleformat{\chapter}[hang]{\filcenter\LARGE\scshape\bfseries}{}{0pt}{} % format nagłówków
\titlespacing*{\section}{0pt}{0pt}{0pt}	
\titleformat{\section}[hang]{\filcenter\large\scshape}{}{0pt}{} % format nagłówków

% ----------- RUBRYKA -------------------
\newcommand{\rubryka}[1]{
	\begin{spacing}{0.8}
	\begin{center}
		\parbox{0.9\textwidth}{\centering\textcolor{red}{\small\textit{#1}}}
	\end{center}
	\end{spacing}
	\linespace}
% ----------- SEKCJA --------------------
% -----------------------------------------
%\renewcommand\fbox{\fcolorbox{red}{white}}% czerwona ramka fbox-a
\newcommand{\sekcja}[2]{
	\begin{spacing}{0.6}
	\begin{center}
		\section*{#1}
		\linespace
		\begin{minipage}{0.88\linewidth}
			\red{\small\textit{#2}}
		\end{minipage}
	\end{center}	
	\end{spacing}
	\paragraphspace}
%------------DEKOR ---------------------
\newcommand{\dekor}[2]{
	\begin{center}
		\vspace*{\fill}
		\includegraphics[width=#1\textwidth]{#2}
		\vspace*{\fill}
	\end{center}}
%------- Zarządzanie wersami w psalmach z tlumaczeniem -------------------------
\newcounter{licznik}
\newcommand{\licznikZero}{\setcounter{licznik}{0}}% zeruj licznik
\newcommand{\licznikUp}{\refstepcounter{licznik}}% zwiększ licznik 
\licznikZero
\newcommand{\BeginParallel}{\begin{spacing}{0.9} \begin{paracol}{2} \licznikZero}
\newcommand{\lat}[1]{\selectlanguage{latin}\ensurevspace{2\baselineskip}\begin{minipage}[t]{\linewidth}\ifnum \value{licznik}>0 \licznikUp \thelicznik .~{#1}\par \else \licznikUp{#1}\par \fi \end{minipage}\switchcolumn}
\newcommand{\latN}[1]{\selectlanguage{latin}\ensurevspace{2\baselineskip}\begin{minipage}[t]{\linewidth}{#1}\par\end{minipage}\switchcolumn}% kolumna bez numeracji wersów
\newcommand{\pol}[1]{\selectlanguage{polish}\begin{minipage}[t]{\linewidth}{#1}\par\end{minipage}\switchcolumn*}
\newcommand{\EndParallel}{\end{paracol} \end{spacing}}

% --- ZNAKI SPECJALNE ---
% 
%\newcommand <c>†</c>
% --------------------------

\begin{document}
\pagestyle{plain}	
\large

\setlength{\parindent}{0pt}
\gresetbracerendering{font}

\pagenumbering{gobble} % nie jest wyprowadzany żaden numer strony, chociaż numer jest nadal resetowany. Odniesienia do tej strony również są puste.
\sekcja{Tłumaczenia\footnotemark}{}
%---------------------
\small
\vspace{-10mm}
\red{\textbf{Incipit}}\\
\red{\Vbar.}~Boże, + wejrzyj ku wspomożeniu memu. \red{\Rbar.}~Panie, pośpiesz ku ratunkowi memu.\\
Chwała Ojcu i Synowi... Chwała Tobie Panie, Królu wiecznej chwały.\\
\red{\textbf{Ant.1}} A Jakób * zrodził Józefa, męża Marji, z której się narodził Jezus, którego zowią Chrystusem.\\
\red{\textbf{Psalm 109}}\\
1. Rzekł Pan Panu memu: * Siądź po prawicy mojej:\\
2. Aż położę nieprzyjaciół twoich * podnóżkiem nóg twoich.\\
3. Berło mocy twojej wypuści Pan z Syjonu: * panuj wpośród nieprzyjaciół twoich.\\
4. Przy tobie panowanie w dzień mocy twej w blaskach świętości: * z żywota przed jutrzenką zrodziłem cię.\\
5. Przysiągł Pan i nic będzie żałował: * Ty jesteś kapłanem na wieki, według porządku Melchizedecha.\\
6. Pan po prawicy twojej, * poraził w dzień gniewu swego królów.\\
7. Będzie sądził narody, dopełni zagłady, * potłucze głowy wielu na ziemi.\\
8. Ze strumienia na drodze pić będzie: * dlatego wywyższy głowę.\\
%---------------------
\red{\textbf{Ant.2}} Posłan jest Anioł * Gabrjel od Boga do panny poślubionej mężowi, któremu było Józef, z domu Dawidowego: a imię panny Marya.\\
\red{\textbf{Psalm 110}}\\
1. Będę Cię wysławiał, Panie wszystkim sercem moim, * w radzie sprawiedliwych i w zgromadzeniu.\\
2. Wielkie są dzieła Pańskie, * doświadczone według wszystkiego upodobania Jego.\\
3. Chwalebne i wspaniałe jest dzieło Jego, * a sprawiedliwość Jego trwa na wieki wieków.\\
4. Uczynił pamiątkę dziwów swoich miłosierny i litościwy Pan: * dał pokarm tym, którzy się Go boją.\\
5. Pamiętać będzie wiecznie na przymierze swoje: * moc uczynków swoich oznajmi ludowi swemu:\\
6. Aby im dać dziedzictwo narodów: * dzieła rąk Jego są prawdą i sądem.\\
7. Pewne są wszystkie przykazania Jego, utwierdzone na wieki wieków, * uczynione w prawdzie i w prawości.\\
8. Zesłał odkupienie ludowi swemu, * ustanowił na wieki przymierze swoje.\\
9. Święte i straszne jest imię Jego: * początkiem mądrości jest bojaźń Pańska.\\
10. Zrozumienie dobre mają, wszyscy, którzy tak czynią: * chwała Jego trwa na wieki wieków.\\
%---------------------
\red{\textbf{Ant.3}} Gdy była poślubiona * Matka Jezusa, Marja Józefowi, pierwej niźli się zeszli, znaleziona jest w żywocie mająca z Ducha Świętego.\\
\red{\textbf{Psalm 111}}\\
1. Błogosławiony mąż, który się boi Pana: * w przykazaniach Jego będzie rozmiłowany wielce.\\
2. Możne na ziemi będzie potomstwo jego: * ród prawych będzie błogosławiony.\\
3. Sława i bogactwo w domu jego, * a sprawiedliwość jego trwa na wieki wieków.\\
4. Wzeszła w ciemnościach światłość prawym, * miłościwy, miłosierny i sprawiedliwy.\\
5. Miły jest człowiek, który się lituje i pożycza, będzie miarkował mowy swe z rozsądkiem, * bo na wieki nie będzie poruszony.\\
6. W wiecznej pamięci będzie sprawiedliwy: * nie będzie się bał pogłosek złych.\\
7. Gotowe serce jego nadzieję mieć w Panu, umocnione jest serce jego: * nie poruszy się, aż wzgardzi nieprzyjaciółmi swymi.\\
8. Rozproszył, dał ubogim, sprawiedliwość jego trwa na wieki wieków, * róg jego podniesie się w chwale.\\
9. Grzesznik ujrzy i będzie się gniewał, będzie zgrzytał zębami swymi i będzie schnął, * żądza niezbożnych zaginie.\\
%---------------------
\red{\textbf{Ant.4}}  Józef, mąż jej, * będąc sprawiedliwym i nie chcąc jej osławiać, chciał ją potajemnie opuścić.\\
\red{\textbf{Psalm 112}}\\
1. Chwalcie, słudzy, Pana, * chwalcie imię Pańskie.\\
2. (ukłonić się) Niech będzie imię Pańskie błogosławione * odtąd i aż na wieki.\\
3. Od wschodu słońca aż do zachodu, * chwalebne jest imię Pańskie.\\
4. Wywyższony jest Pan nad wszystkie narody * i nad niebiosa chwała Jego.\\
5. Któż jest jak Pan, Bóg nasz, który mieszka na wysokości, * a na niskość spogląda na niebie i na ziemi?\\
6. Podnosi z ziemi biednego * i z gnoju wywyższa ubogiego:\\
7. Aby go posadzić z książętami, * z książętami ludu swego.\\
8. Który czyni, iż niepłodna mieszka w domu, * matka synów z weselem.\\
%---------------------
\red{\textbf{Ant.5}} Anioł Pański * ukazał się Józefowi mówiąc: Józefie, Synu Dawidów, nie bój się przyjąć Marji małżonki twej: albowiem co się w niej narodziło jest z Ducha Świętego: a porodzi syna i nazwiesz imię Jego Jezus.\\
\red{\textbf{Psalm 116}}\\
1. Chwalcie Pana wszystkie narody, * chwalcie Go wszystkie ludy:\\
2. Bo możne jest nad nami miłosierdzie Jego, * a prawda Pańska trwa na wieki.\\
%---------------------
\red{\textbf{Capitulum}}\\
Mąż wierny wielce chwalon będzie. A który strzeże Pana swego, uwielbiony. (Prz 28,20; 27,18)\\
%---------------------
\red{\textbf{Hymn}}\\
1. Ciebie, Józefie, sławią niebieskie zastępy, * Niechaj Twą chwałę głoszą także chrześcijan chóry * Co zasługami jaśniejący, czystym związkiem * Złączony jesteś z Najchwalebniejszą Dziewicą.\\
2. Zdumiony widzisz, że Twa Oblubienica ma zostać Matką. * Zwątpienie Cię przenika. * Anioł dopiero w widzeniu Ci objawia, * Że to Poczęcie tchnieniem Płomienia z wysoka.\\
3. Narodzonego Pana w objęcia przyjmujesz, * Uchodzisz z Nim w daleką krainę Egiptu, * W Jerozolimie zgubionego szukasz, płacząc, * I znajdujesz mieszając radość swą ze łzami.\\
4. Wszakże po zgonie wieczna szczęśliwość * Udziałem zwykłych śmiertelnych, śmierć palmę chwały im wręcza: * Ty już za życia, jak Święci, zażywasz Boga, * Przedziwnym losem jesteś tu od nich szczęśliwszy.\\
5. Najwyższa Trójco, przepuść kornie błagającym. * Daj dla zasług Józefa, do nieba się dostać. * Aby nam kiedyś, sługom Twym, wolno było, * Nieustający Tobie hymn wdzięczności śpiewać. * Amen.\\
%---------------------
\red{\Vbar.}~Postanowił go panem domu swego.\\
\red{\Rbar.}~I książęciem wszystkiej posiadłości swojej.\\
%---------------------
\red{\textbf{Ant.Magnificat}}~~A Józef, wstawszy * ze snu, uczynił jako mu rozkazał anioł; i przyjął żonę swoją.\\
\red{\textbf{Magnificat}}\\
1. Wielbij * duszo moja, Pana.\\
2. I rozradował się duch mój: * w Bogu, Zbawicielu moim.\\
3. Iż wejrzał na niskość służebnice swojéj: * albowiem oto odtąd błogosławioną mię zwać będą wszystkie narody.\\
4. Albowiem uczynił mi wielkie rzeczy, który możny jest: * i święte imię jego.\\
5. A miłosierdzie jego od narodu do narodów: * bojącym się jego.\\
6. Uczynił moc ramieniem swojem: * rozproszył pyszne myśli serca ich.\\
7. Złożył mocarze z stolice: * a podwyższył niskie.\\
8. Łaknące napełnił dobrami: * a bogacze z niczem puścił.\\
9. Przyjął Izraela, sługę swego: * wspomniawszy na miłosierdzie swoje.\\
10 Jako mówił do ojców naszych: * Abrahamowi i nasieniu jego na wieki.\\
%---------------------
\red{\textbf{Oratio}}\\
Módlmy się. Niechaj zasługi Oblubieńca Przenajświętszej Rodzicielki Twojej wspomogą, nas, prosimy Cię Panie, aby to, czego własnemi siłami otrzymać nie możemy, za jego wstawiennictwem było nam dane. Który żyjesz i królujesz z Bogiem Ojcem w~jedności Ducha Świętego Bóg przez wszystkie wieki wieków. Amen.\\
\red{\textbf{Commemoratio}}\\
\red{\textbf{Ant.}}~~Albowiem jeden jest nauczyciel wasz, a wszyscy jesteście bracia. I ojca nie zówcie sobie na ziemi; albowiem jeden jest Ojciec wasz, który jest w niebiesiech. Ani się zówcie nauczycielmi, gdyż jeden nauczyciel wasz, Chrystus. (Mt 23,8-10)\\
\red{\Vbar.}~Aniołom swoim Bóg rozkazał o tobie. (Ps 90,11-12)\\
\red{\Rbar.}~Aby cię strzegli na wszystkich drogach twoich. \\ 
Módlmy się. Daj się przejednać, Panie, naszymi modłami i ulecz choroby dusz naszych, abyśmy otrzymawszy przebaczenie, zawsze się cieszyli Twoim błogosławieństwem. Przez naszego Pana Jezusa Chrystusa, Twojego Syna, który z Tobą żyje i króluje w jedności Ducha Świętego, Bóg, przez wszystkie wieki wieków.\\
\red{\textbf{Conclusio}}\\
\red{\Vbar.}~Panie, wysłuchaj modlitwy mojej. \red{\Rbar.}~A wołanie moje niech do Ciebie przyjdzie.\\
\red{\Vbar.}~Błogosławmy Panu. \red{\Rbar.}~Bogu niech będą dzięki.\\
\red{\Vbar.}~Dusze wiernych zmarłych przez miłosierdzie Boże niech odpoczywają w pokoju. \red{\Rbar.}~Amen.\\
Ojcze nasz, któryś jest w niebie, święć się imię Twoje: przyjdź królestwo Twoje: bądź wola Twoja, jako w niebie tak i na ziemi. Chleba naszego powszedniego daj nam dzisiaj: i odpuść nam nasze winy, jako i my odpuszczamy naszym winowajcom: i nie wódź nas na pokuszenie: ale nas zbaw ode złego. Amen.\\
\red{\Vbar.}~Niech Pan nas obdarzy swoim pokojem. \red{\Rbar.}~I życiem wiecznym. Amen.\\
\red{\textbf{Antiphona finalis B.M.V.}}\\
Witaj, niebios Królowo, * Witaj, Pani aniołów, * Witaj, Różdżko, witaj Bramo, * Jasność zrodziłaś światu.\\
Ciesz się, Panno chwalebna, * Ponad wszystkie piękniejsza, * Witaj, o Najśliczniejsza, * Proś Chrystusa za nami.\\
\red{\Vbar.}~Dozwól mi chwalić Cię, o Pani święta. \red{\Rbar.}~Daj mi moc przeciw nieprzyjaciołom Twoim.\\
Módlmy się. Udziel nam, prosimy Cię, miłosierny Boże, pomocy w słabościach naszych, abyśmy, którzy świętej Bożej Rodzicielce cześć oddajemy, za jej przyczyną i wstawiennictwem, powstali z nieprawości naszych. Przez tegoż Chrystusa, Pana naszego. Amen.\\
\red{\Vbar.}~Pomoc Boża + niech będzie zawsze z nami. \red{\Rbar.}~Amen.\\
\footnotetext{Tłumaczenia pochodzą ze strony divinumofficium.com oraz z Biblii w tł. ks. Jakuba Wujka SJ}
\end{document}